\chapter{Instruction encoding}\label{ch:instruction-encoding}

A G17 program is a byte stream of variable-length instructions aligned only to two bytes. Apple's toolchain has no
mnemonics for them. An operand's bits are spread across the instruction, often non-contiguous, sometimes inverted, and
placed differently at each length an opcode takes. Apple ships a decoder for this stream inside its GPU compiler but
withholds the printer. The project calls the decoder directly, and every structural statement in this reference comes
from it; statements about execution come from dispatches.

Three kinds of statement about an encoding are easily confused: \emph{what Apple's decoder recognises} (an opcode, a
length, typed operands), \emph{what can be authored} (bytes this project's encoder can produce for a requested operand
set, read back through the decoder) and \emph{what the hardware does with those bytes}.

Errors in this project have come from each boundary between them. They include a decoder-clean back edge that was wrong in its
whole second half, a "free" bit that changed a compare's result, a refused tensor bit that changed the product, an
encoder that rounded an odd register silently, and a name inferred from a structural twin that was wrong 44 times in 44.
Each rule below says which of the three it concerns, and \cref{sec:admission-decodes-executes} gives the guards that
check authored bytes against each boundary.

Terms:

\begin{itemize}
\item \emph{Form}: an (opcode, length) pair, written for example \op{12674}/16. Field maps are per form, never per opcode (\cref{sec:framing}).
\item \emph{byteN{[}k{]}}: bit k (0 = least significant) of byte N of the instruction, byte 0 first in the stream. \emph{Bit n} (linear):
  byte n // 8, bit n \% 8.
\item \emph{Operand}: a position in the decoder's operand list (operand 0 first). Registers print as Rn (\cref{ch:register-file}); immediates as
  integers; a base address as \code{expr:} (a pointer into the decoder's process, meaningless outside it).
\item \emph{Operand 1}: in ALU, memory and tensor forms, the destination's raw operand, printed as one wide integer (older notes:
  the "flags word" or "control word"). It packs the destination modifier (bits 4–5), a fill slot (bits 20–23) and a wait
  mask (bits 24–31) (\cref{sec:scoreboard-publish-wait}).
\item \emph{Witness}: an instance of a form in Apple's compiled code. \emph{Repair-walk witness}: bytes reached by mutating a
  neighbour until the decoder accepts them, with no Apple instance.
\item \emph{D1, D2, D3}: the three coverage counts (\cref{sec:opcode-identity-counts}).
\item \emph{Corpora}: \code{isa/g17-corpus-programs.jsonl} (6,594 programs, 365,590 instructions) and the vendor corpus (9,556
  functions, 4,644,438 instructions), union 1,893 forms (MM~\mm{25.51}, MM~\mm{25.68}, MM~\mm{25.78}).
\end{itemize}

\section{Instruction lengths and framing}\label{sec:framing}

Lengths are even, from 2 to 22 bytes, and belong to the instance, not the opcode (decoder). Apple's decoder frames
three populations as follows.

\begin{itemize}
\item \emph{Apple's two corpora} (5,010,028 instructions, 1,893 forms; \code{isa/g17-vendor-corpus-forms.json}): every even length
  from 2 to 16, plus 18 bytes (4 forms, 31,458 instructions, among them \op{10909}) and 22 bytes (30 forms,
  53,629 instructions, among them \op{14661}); no 20-byte instance. The 18- and 22-byte instances are
  1.7 percent of the total.
\item \emph{The 216-opcode census} (\code{isa/g17-opcodes.toml}, \code{docs/agx3-oracle.md} section 7) spans 2 to 16 bytes; none of
  the 18- or 22-byte opcodes above is in it, so a range read from this census stops at 16.
\item \emph{The decoder-admitted universe} (\val{admitted-opcodes} opcodes, \cref{sec:opcode-identity-counts}): the 6,000 opcodes with no Apple witness take their width
  from decoding a repair-walk encoding, which gives 4 to 16, 18 and 22 bytes (7 opcodes at 18, 38 at 22) and again no
  20 (\code{tools/g17isamap.py} \code{repair\_walk\_widths}). Such a width means the decoder accepts the form, not that anything
  emits it. Whether a 20-byte length exists is not known; no source contains one and no source excludes it.
\end{itemize}

In the 216-opcode census, 64 opcodes occur at more than one length: \op{12674} at 12 or 16 bytes,
\op{554} at 4 or 8, \op{11842} at 2, 8 or 10
(\code{docs/agx3-oracle.md} section 7). Field positions move with the length: \op{2190}'s operand-1
high code sits at $2^{24}$-$2^{28}$ in its 8-byte form, $2^{30}$-$2^{31}$ in the 16-byte form and $2^{32}$/$2^{36}$ in the 12-byte form (MM~\mm{25.53}); one byte
offset can mean different fields in loads and stores of different lengths (MM~\mm{25.24}), so field maps must be per form.
One map per opcode was the common defect behind five separate fixes: for 1,301 of 1,401 per-form specifications the
opcode-wide authoring map had been probed on an instance of another length (\op{13575}'s 16-byte map reached
byte 9 of its 4-byte form) (MM~\mm{25.88}).

\Cref{tab:length-rules} collects the known length rules.

\begin{xltabular}{\linewidth}{@{}X>{\hsize=1.18\hsize}X>{\hsize=0.85\hsize}X@{}}
\caption{Known length rules by population, each a statement about Apple's decoder's framing.}\label{tab:length-rules}\\
\toprule
Population & Rule & Standing \\
\midrule
\endfirsthead
\tablecontinued{3}
\toprule
Population & Rule & Standing \\
\midrule
\endhead
general & lengths are a base plus optional extensions; the length-bit cluster
\{32, 63, 64, 79, 80\} is conserved from G13 & decoder, 895-opcode \code{length\_bits} work (MM~\mm{25}) \\
class 7 (byte0 low nibble 7), byte1 = 0, byte6 in \{0x60, 0x68, 0x70,
0x78\} & length = 12 + 4 × byte10{[}0{]} & decoder; 9,913 instances (MM~\mm{25}) \\
class 7, byte1 = 0, byte6 in \{0x83, 0xa3, 0xa4, 0xa9, 0xaa, 0xab, 0xae,
0xaf\} & constant per byte6 (12 or 10) & decoder; 13,571 instances (MM~\mm{25}, \code{agxforge/g17/dis.py}) \\
class 7, byte1 = 0, byte6 = 0x10 / 0xa0 / 0xa2 & bit 63 / bit 18 / bit 18 selects the length & causal against the decoder, 60 of 60, 37 of 37, 15 of 15 decodable flips
(MM~\mm{25}) \\
class 7, byte1 = 0, byte6 = 0xa1 (lengths 10, 12, 14, 16) & open; bit 45 explained 81.5 percent correlationally and moved the length
in 0 of 60 flips & bounded negative (MM~\mm{25}) \\
\op{11842} & 2 bytes exactly when bit 15 = 0 (1,682 of 1,682); the 10-byte driver is
not a function of the always-present bits & decoder (MM~\mm{25.62}) \\
\op{17235}, \op{17244} & 8 bytes exactly when bit 63 = 0; 14 needs bit 63 = 1, not sufficient & decoder (MM~\mm{25.62}) \\
\op{12709} & 14-byte form when it carries bit 37 or a displacement above 255 & decoder and Apple's compiler (MM~\mm{6}) \\
\op{592} & 4-byte form's K field stops at 126; Apple uses the
8-byte form for K 128–134 & Apple's compiler emits (MM~\mm{25.120}) \\
every opcode & bits 0, 1, 2 never vary within an opcode & decoder (MM~\mm{25.62}) \\
\bottomrule
\end{xltabular}

For class 7, \code{docs/agx-isa-status.md} states, from a tiling solution, that 16 bytes need \emph{both} byte10
bit 7 and bit 0 (0x21 → 12, 0x80 → 12, 0x81 and 0xc1 → 16); MM~\mm{25} quotes \code{\_WIDE7\_EXT} as 12 + 4 × byte10{[}0{]}, which
reads 0x21 as 16. They were derived from different populations. Checked against the lengths Apple's decoder assigned in the corpus's span
index: in the \code{\_WIDE7\_EXT} population (byte6 in \{0x60, 0x68, 0x70, 0x78\}) byte10 takes only 0x80, 0x81, 0xc1 and once
0xe1, bit 7 is always set, and 12 + 4 × byte10{[}0{]} holds on all 12,001 such class-7 spans (9,913 of them with byte1 = 0,
MM~\mm{25}'s count); byte10 = 0x21 occurs only in other byte6 groups (0xa3, 0xa2, 0x83, 0xa8, 0xa0), where its instances are
12 bytes except under 0xa8 (16) and 0xa0 (14). Each formula holds on its own population and the two agree where they
overlap, but neither is a general class-7 rule: "both bits" fails for byte6 = 0xa8 (0x21, 16 bytes). Tiling alone cannot settle the question,
because two disjoint length assignments satisfy all 364 tiling constraints (\code{docs/agx-isa-status.md}). The running
example follows the \code{\_WIDE7\_EXT} rule: its tensor loads are 12 bytes with byte10 = 0x80 and 16 bytes with byte10 = 0x81, and in that
listing the 16-byte ones are exactly those with a displacement of 256 bytes or more.

Ground truth for framing is Apple's decoder (\cref{sec:decoder}); the project's own framer cannot replace it. The framer "landed" exactly on
22,184 of 22,187 constant programs but framed them identically to Apple in only 1,725 (7.8 percent), because a wrong
framing falls back into phase on the short fixed tail of end-of-program and two-byte fillers; it misframes 200 of 200
current programs within their first instructions (MM~\mm{25.88}, MM~\mm{25.209}).

\section{Opcode identity and coverage counts}\label{sec:opcode-identity-counts}

Opcode numbers are indices into Apple's instruction table as built for macOS 26.6.2 (25G83). That MCInstrDesc table has
17,796 entries (stride 48; the first halfword equals the index for all of them), of which 17,779 declare at least one
operand (\code{isa/g17-agx3meta-instrs.txt}; sources quoting one or the other are reconciled in \cref{app:a}). The macOS 27
build's table has \val{renumber-new-opcodes} entries in a different order; this reference keeps the 25G83 numbering
(\cref{sec:decoder-renumbering}). Apple built its LLVM with instruction names disabled, so the name table holds
"0", "1", "2" and there is no assembly-string table; every opcode name in this reference is a project name from measured
function, recorded in \code{isa/g17-opcode-glossary.json}, not an Apple spelling (decoder, \code{docs/agx3-oracle.md} section 6).
Mechanically available per opcode: operand count, number of definitions, register class and operand type per operand
(type 4 = PC-relative), scheduling class, flag words; lengths are not a table property (\code{agxforge/g17/model.py}).

The decoder admits \val{admitted-opcodes} opcode records (6,718 found by mutation and breadth-first walks
from real encodings, plus \op{13754}, added from the vendor corpus) (MM~\mm{24.8}
row T1, MM~\mm{25.58}). The declared set overstates the ISA: in the atomic family 208 opcodes are declared and 29 admitted
(\code{tools/g17admit.py}).

\Cref{tab:form-counts} gives the three form counts (generated, \code{docs/g17-isa-reference.md}); they are never summed.

\begin{xltabular}{\linewidth}{@{}>{\hsize=0.75\hsize}Xl>{\hsize=1.25\hsize}X@{}}
\caption{The form counts D1, D2 and D3, with the dispatched forms outside the verified class.}\label{tab:form-counts}\\
\toprule
Count & Forms & What it requires \\
\midrule
\endfirsthead
\tablecontinued{3}
\toprule
Count & Forms & What it requires \\
\midrule
\endhead
D1 structural & \val{d1} & the decoder can address the form (opcode, length); 7,549 with the 25G83 decoder
(\cref{sec:decoder-renumbering}) \\*
D2 encoded & \val{d2} & an encoding Apple wrote, with no unforced and no undetermined bits \\*
D3 own instruments (any of nine) & \val{d3-any} & an own hardware instrument checked its semantics (D3 verified class, an
isolated one-opcode probe returning the predicted value: \val{d3-verified}) \\*
dispatched, not in the verified class & \val{dispatched-unverified} & it ran; no isolated probe confirmed the name (\val{dispatched-other-evidence} of these carry other
semantic evidence) \\*
dispatched, no semantic evidence of any kind & \val{dispatched-no-evidence} & it ran and nothing has read its output \\*
\bottomrule
\end{xltabular}

\begin{itemize}
\item D1 is known to undercount. Five tensor MMA forms (fp32 × fp32 and 16-bit A × fp32 B, each with and without
  accumulator, and int8 widening without accumulator) are admitted by the decoder but missing from the admission walk:
  a selector bit that also changes an operand's register class cannot be reached one flip at a time
  (\code{docs/g17-isa-reference.md}). D1 counts what the contract records, not what the decoder accepts (MM~\mm{24.8} row T1).
\item D2 counts forms fully specified by Apple-witnessed, bit-settled evidence, which is a different property from being
  emitted by the compiler; 91.7 percent of D2 forms reconstruct their own witness and 57 do not (MM~\mm{25.87}). D2 is held back by missing per-form evidence:
  6,579 of 6,858 non-D2 forms failed on absence, not on unknown bits. Its last gains came from hardware checks of
  drifting bits (MM~\mm{25.87}, MM~\mm{25.58}; the successive counts are in \cref{app:a}).
\item The nine D3 columns are instruments with different failure modes (a name an isolated probe confirmed; the one library
  candidate surviving a form's cases; a family law), so only "any of the nine" is quoted (\code{docs/g17-isa-reference.md}).
\item \emph{Correction:} \code{docs/g17-isa-cartography.md}'s D1 975, D2 690 and two-MMA admission are a superseded snapshot
  (\cref{app:a}).
\end{itemize}

The counts do not give the number of independent semantic questions. The earlier "527" is retracted (21 of 25
multi-member name-stem classes contain opcodes with different functions), and naming by structural twin was right 0
times in 44 on the verified set (MM~\mm{25.81}, \code{docs/g17-isa-cartography.md}). Forms Apple's compiler never emits
are absent from both corpora, so corpus reading cannot find them (MM~\mm{25.78}).

\section{Field maps}\label{sec:field-maps}

A field map is linear per form: \code{decoded operand = base + sum(bit\_is\_set[i] x weight[i])}, negative
weights for inverted bits (MM~\mm{23.3}). Operand domains (\code{agxforge/g17/auth.py}): \emph{raw} (the value the decoder prints:
immediates, modifier words), \emph{slot} (register x 2 + half, ordinary registers), \emph{index} (class-relative member index, for
aligned tuple classes; \cref{sec:register-classes}). Each instruction bit is classified by flipping it and asking the decoder
(\code{agxforge/g17/encode.py}): \emph{opcode} (the flip decodes as another opcode), \emph{forced} (the flip is undecodable; the recorded
value is the only legal one), \emph{operand} (the flip moves a printed operand), \emph{invisible} (nothing printed changes).
The classification has four complications.

\begin{itemize}
\item A value bit can have several carriers that are not interchangeable. \op{3290}'s operand 3 value bit 4 is carried by
  byte6{[}3{]} inverted and by byte8{[}2{]}; byte6{[}3{]} also turns operand 2 into an expression. 1,031 of 9,282 sampled operands
  have such a bit (\code{agxforge/g17/auth.py}).
\item A bit's role can depend on other bits: 547 of 1,401 per-form specs carry an opcode- or forced-class bit that Apple's
  own instances vary (1,132 bits; MM~\mm{25.87}, MM~\mm{25.88}).
\item Some immediates are table codes rather than literal values (\op{13754}'s operand 2 prints 8, 224,776 and 183,494); the codec refuses to write them as numbers (MM~\mm{25.58}).
\item An invisible bit can still change what executes (\cref{sec:admission-decodes-executes}).
\end{itemize}

A field map counts as established when it rebuilds Apple's instances. In the best case, a 38-bit map for \op{11462}/10, fitted on 1,637 instances, rebuilds 33,473 of 33,475 vendor instances with zero
differences (decoder, MM~\mm{25.52}). Across the corpus the rates are lower: corpus-fitted maps re-author 92.0 percent of in-sample
instructions byte-exactly but 55.4 percent of held-out vendor instructions (2026–09-06 snapshot,
\code{docs/agx-isa-status.md}); the encoder's held-out re-authoring now stands at 80.8 percent (MM~\mm{25.88}).

\textbf{Worked example: the operand fields of a tensor MMA} (compiled, decoded; executed in the running example). The first MMA of
the running example (offset 0x29a), \op{5107}, is
\code{a7 04 05 11 22 04 a4 62 02 02}. Apple's decoder prints
\code{R24\_...\_R31 <- 16777216, 1, R20\_R21\_R22\_R23, 0, 2, R68\_R69\_R70\_R71, 0, 2}. The field map
(\code{isa/g17-tensor-mma-fieldmap.json}, entry 5107, reference \hash{2f0005102200a4420200}) reproduces every operand from the
bytes; it records each bit's effect as a flip from its reference, and this instance differs from the reference in seven
bits (3, 7, 10, 24, 42, 61, 73). \Cref{tab:mma-operand-fields} lists each operand and its carriers.

\begin{xltabular}{\linewidth}{@{}l>{\hsize=0.44\hsize}X>{\hsize=1.35\hsize}X>{\hsize=1.21\hsize}X@{}}
\caption{Operands of the running example's first tensor MMA and the bits that carry them in this instance.}\label{tab:mma-operand-fields}\\
\toprule
Operand & Value & Carried by (in this instance) & Meaning \\
\midrule
\endfirsthead
\tablecontinued{4}
\toprule
Operand & Value & Carried by (in this instance) & Meaning \\
\midrule
\endhead
0, D & R24 & byte0{[}7{]} = 1 (+8), byte7{[}5{]} = 1 (+16); byte2{[}3{]},
byte2{[}4{]}, byte2{[}7{]} = 0 & 8-register destination; field in units of 8 registers \\*
1 & 16,777,216 = $2^{24}$ & byte1{[}2{]} = 1 ($2^{24}$); byte0{[}3{]} = 0 (no $2^{31}$); byte4{[}3{]}
= 0 (no keep, weight 32) & wait on scoreboard slot 0 (the tile loads publish slot 0); not a keep,
since the next MMA writes another tuple \\*
2 & 1 & no variable bit in the map & a constant of this form \\*
3, A & R20 & byte3{[}0{]} = 1 (+16 from the reference's R4) & 4-register tuple \\*
4 & 0 & byte1{[}7{]} (32), byte2{[}5{]} (16), both 0 & A's modifier: neither keep nor release (
\cref{sec:release-does}) \\*
5 & 2 & byte6{[}2{]} = 1 (clear gives 3); byte8{[}2{]} = 0 (transpose adds 32) & A is half (3 would be bfloat) \\*
6, B & R68 & byte5{[}2{]} = 1 (+64), byte9{[}1{]} = 1 (+4) & 4-register tuple \\*
7 & 0 & byte5{[}1{]} (16), byte8{[}0{]} (32) & B's modifier \\*
8 & 2 & byte7{[}6{]} = 1; byte8{[}3{]} = 0 (transpose) & B is half \\*
\bottomrule
\end{xltabular}

The remaining bits are structural or inert: byte8{[}1{]} = 1 selects this no-accumulator form (clearing it decodes as
\op{5106}); byte5{[}7{]}, byte8{[}4{]} and byte6{[}0{]} would select fp32-B, fp32-A and the
int8 family (MM~\mm{0.4}, MM~\mm{2}); bits 0–2 are constant. Across the accumulate form's reference all 80 bits are accounted for:
39 claimed by mapped fields, 25 structural, 16 decoder-inert, none unaccounted (MM~\mm{6}). Its hardware single-bit census ran
80 of 80 flips with none refused by the pipeline and none faulting; every decoder-inert bit was hardware-inert (measured,
MM~\mm{2}). The project's MMA encoder reproduces 800 of 800 compiled MMAs byte-exactly from decoded operands (MM~\mm{2}).

\textbf{Worked example: a tensor load's operand 1} (compiled, decoded). The first tile load,
\code{47 00 47 0a 08 00 60 2a 00 80 80 08}, is \op{12674}/12:
\code{R20\_R21 <- 17825792, 241, expr, 0, R5, 0, 0, 2, 15}: destination pair (operand 0), operand 1, the base address pair
(operand 3, printed as an expression), index R5 with modifier 0 (operands 5, 6), displacement 0 (operand 7), element code
2 (operand 8) and mask 15 (operand 9); operands 2 and 4 are immediates not discussed here. Operand 1 is 17,825,792 = 0x01100000,
exactly what \code{agxforge/g17/memenc.py} authors as \code{((slot + 1) << 20) | (wait\_mask << 24)} with slot 0 and wait mask
0b00000001:

\begin{itemize}
\item bits 20–23 = 1: this load \emph{publishes} scoreboard slot 0 (the decoder prints slot + 1; 0 means none);
\item bit 24: this load \emph{waits} on slot 0, where the special-register read that feeds its index R5 published.
  Single-bit flips of this instance through Apple's decoder locate the carriers: clearing byte7{[}5{]} removes $2^{24}$ (the wait);
  setting byte6{[}3{]} or byte6{[}4{]} changes the slot field (to 2 and 3); setting byte2{[}5{]} adds 32, the destination keep weight
  (decoder only, this instance). The slot carriers differ between the tensor load's four field variants (12 and 16 bytes, with
  and without a residue bit byte1{[}6{]}), so they must not be copied from one variant to another (MM~\mm{25.26}, MM~\mm{25.32}). The
  element code (operand 8) is the access width in bytes, 1, 2, 4, 8 or 16 across the load family (MM~\mm{25.37}). The
  index's modifier is 0 because the next load reads R5 again; a load releases its index with 16 (\code{memenc} \code{index\_last}).
\end{itemize}

\textbf{Worked example: the scalar tail} (compiled, decoded, executed). The running
example's last five instructions compute \code{C[tg.x] = A'[tg.x] + 1.0} on a scalar path (\cref{tab:example-scalar-tail}).

\begin{xltabular}{\linewidth}{@{}l>{\hsize=0.57\hsize}XX>{\hsize=1.36\hsize}X@{}}
\caption{The running example's scalar tail, five instructions from special-register read to store.}\label{tab:example-scalar-tail}\\
\toprule
Offset & Bytes & Decoded & Fields that matter \\
\midrule
\endfirsthead
\tablecontinued{4}
\toprule
Offset & Bytes & Decoded & Fields that matter \\
\midrule
\endhead
0x49e & \code{9c9c1006} & \op{14059}: R9 ← 1048576, \code{SR\_TG\_X} & operand 1 = $2^{20}$: publishes slot 0. Destination R9 in byte0{[}7:4{]}
(a four-bit field, \cref{sec:narrow-fields-narrow}); special
register number in byte1 \\*
0x4a2 & \hash{0f084312182210c0410080000000} & \op{12682}/14: R16 ← 137464119296, 2066,
expr, 0, R9, 0, 0, 4 & operand 1 = $2^{37}$ + 8 × $2^{20}$ + $2^{24}$: publishes slot 7, waits on
slot 0 (for R9); index R9 kept (read again by the store); element code 4 \\*
0x4b0 & \code{1c80423e6020000c} & \op{11842}/8: R17 ←
68736253952, 1065353216 & immediate 0x3F800000 = 1.0f, scattered (below) \\*
0x4b8 & \hash{0900044a0220a02a24081a00} & \op{998}/12: R18 ← 148176371712, R16, 16, R17, 16 & operand 1 = 0x2280000000: bit 31 waits on slot 7 (the load); both
sources released (last reads) \\*
0x4c4 & \code{2f08431201061040} & \op{17229}/8: R18, 16, 2066, expr, 0, R9, 16, 0, 4 & value R18 and index R9 both released; distinct registers, so legal
(\cref{sec:release-does}) \\*
\bottomrule
\end{xltabular}

The move immediate's 32-bit value is scattered across six byte ranges, a layout predicted and confirmed byte-exactly on
0xDEADBEEF and 0x12345678 in the scalar campaign (\code{docs/agx-isa-status.md}): value bits 6:0 in byte1{[}6:0{]}, 10:7 in
byte4{[}4:1{]}, 12:11 in byte5{[}3:2{]}, 20:13 in byte6, 24:21 in byte7{[}3:0{]}, 31:25 in byte3{[}7:1{]}. Here byte3 = 0x3e gives bits
25–29, byte7 = 0x0c gives bits 23–24, everything else is zero: 0x3F800000. Its destination R17 is seven bits: byte0{[}7:4{]}
= 1 plus byte2{[}6{]} (+16); single flips of this instance through the decoder give byte2{[}7{]} = +32 and byte2{[}3{]} = +64 (decoder
only). Operand 1's bits 33 and 37 in the load and the add are not explained here.

\section{Scoreboard field encoding}\label{sec:scoreboard-publish-wait}

Semantics (what a slot counts, when a wait is needed, which producers are late) are in \cref{sec:scoreboard}. A producer of a late value names the slot it fills, printed as slot + 1 in operand 1
bits 20–23; a consumer names the slots it waits on in operand 1 bits 24–31, bit 24 + s for slot s; there is no separate
wait instruction (measured, MM~\mm{6}, MM~\mm{0.8}). \Cref{fig:control-word} draws the layout, and \cref{tab:wait-fields} gives the carriers in each form.

% Figure 3.1: the scoreboard fields of operand 1 (MM 6, MM 0.8).
\begin{figure}[htb]
\centering
\begin{bytefield}[bitwidth=1.15em, bitheight=6ex, endianness=big]{32}
  \bitheader[lsb=0]{0,19,20,23,24,31} \\
  \bitbox{8}{\small wait mask\\[-2pt]{\scriptsize bit $24 + s$ waits on slot $s$}}%
  \bitbox{4}{\small fill\\[-2pt]{\scriptsize slot + 1}}%
  \bitbox{20}[bgcolor=codebg]{\small form-specific fields}
\end{bytefield}
\caption{Scoreboard fields in operand 1, for the forms that carry them there. A producer of a late value writes the
slot it fills as slot + 1 in bits 20--23 (0 means none); a consumer sets bit $24 + s$ for each slot $s$ it waits on.
The \op{12682} of the running example carries \code{0x01800000}, so it waits on slot 0 and fills
slot 7. Other forms place the same fields elsewhere (\cref{tab:wait-fields}).}
\label{fig:control-word}
\end{figure}


\begin{xltabular}{\linewidth}{@{}X>{\hsize=0.83\hsize}X>{\hsize=1.21\hsize}X@{}}
\caption{Where each form encodes its scoreboard fill slot and wait mask. Sources follow the table.}\label{tab:wait-fields}\\
\toprule
Form & Fill slot & Wait mask \\
\midrule
\endfirsthead
\tablecontinued{3}
\toprule
Form & Fill slot & Wait mask \\
\midrule
\endhead
general loads, atomics, \op{14059}, imageblock
load & operand 1 bits 20–23 (slot + 1) & operand 1 bits 24–31 \\
\op{12674} & operand 1 bits 20–23; carriers per field variant (
\cref{sec:field-maps}) & operand 1 bits 24–31 (byte7{[}5{]} carries bit 24 in the 12-byte variant
above) \\
\op{12709} & byte4{[}3{]} + 2 × byte4{[}4{]} + 4 × byte6{[}4{]}, printed as slot + 1 & operand 1 bits 24–31 \\
tensor MMA (same carriers in all ten field maps; causally measured on
the 16-bit forms) & none & slots 0–4 at byte1{[}2..6{]}, slot 5 at byte7{[}7{]}, slot 6 at
byte2{[}6{]}, slot 7 at byte0{[}3{]} \\
12-byte ALU forms & none & byte0{[}3{]} = slot 7; byte1{[}2{]} = slot 0 on \op{10279} and \op{10282} \\
14-byte stores & none & byte9{[}5{]} = slot 7 (the carrier this compiler's
encoder writes) \\
\op{17229}, 8-byte form & none & operand 1 bits 24–31 on 15 corpus instances, each matching its pending
slots; this compiler's encoder has no carrier for them \\
\op{17256} & none & operand 1 bits 24–31, every bit used (13,140 instances in a 314-object
census, each one-hot value present) \\
\op{17257}, \op{17258} & none & zero in all 876 and 118 compiled instances \\
\op{592} & 0 into the uniform file; 1–8 (slot + 1) into the staging file; 8-byte
form: field in byte5{[}0..2{]}, value + 1 & none \\
\bottomrule
\end{xltabular}

Sources: MM~\mm{6} (MMA mask causal over 43 dispatches, then 20 of 20 preregistered; vector-load formula, correcting an
inverted one; the vector store's full mask, refuting an earlier "three-bit store limit" that was an encoder-table
artifact; the tensor stores' zero byte), MM~\mm{25.121} (ALU slot-0 bit), MM~\mm{25.120.1} (\op{592} over
178,582 vendor instances: uniform writes always 0, staging writes 1 in the 4-byte form and 1–8 in the 8-byte form, none
crossing), MM~\mm{25.146} and \code{isa/g17-wait-mask-census.json} (the 8-byte store), \code{agxforge/g17/memenc.py}, \code{agxforge/g17/cc.py}.
The decoder census behind the table (every instance's wait bits equal the slots its sources are still pending on, 126
forms, MM~\mm{25.96}) and what each mask means are \cref{sec:scoreboard}. A form whose wait bits are not located (for example \op{798}) gets a loaded source through a waited copy (MM~\mm{25.196}).

byte4{[}3{]} has three unrelated meanings. It is the first term of the vector load's slot formula, an opcode selector
in \op{12674} (flipping it gives \op{12673}, with 9 operands instead of 10), and the keep bit
in the tensor MMA (MM~\mm{24.2} row R6, MM~\mm{25.1}), so a measurement of it in one of these forms does not carry
over to the other two.

Two earlier misreadings took a regularity in Apple's code for a hardware mechanism.

\begin{itemize}
\item \emph{The store "consumer token".} Early analysis found bits in stores that looked like a token with a wait counter, a
  countdown rule and a loop rule. They are bits 5:4 of the store's source register number, on 632 of 632 stores; the
  "rules" were Apple's compiler allocating registers in descending order (MM~\mm{19}).
\item \emph{The MMA "result ring".} The MMA's high operand-1 bits appeared one-hot, rotating through eight positions, and were read
  as tags of the accelerator's in-flight results. They are the eight-bit wait mask over load slots; they looked one-hot
  because each MMA waited on exactly its own load group (MM~\mm{19}, MM~\mm{6}). The causal test that settled this is in
  \cref{sec:scoreboard}.
\end{itemize}

\section{Program boundaries}\label{sec:program-boundaries}

\begin{itemize}
\item The delivered \code{\_\_text} is a 64-byte prologue, the main program, then filler, in 421 of 421
  delivered archives; the prologue in the released class is \op{684} followed by thirty \op{13483}, \code{0e000000} then thirty \code{0600}, and the entry offset is 64 (this compiler and the ABI-v5 class, MM~\mm{23.2},
  MM~\mm{25.148}). The scalar campaign stripped the same fixed 0x40-byte prologue from Apple's minimal kernels
  (\code{docs/agx-isa-status.md}).
\item \op{684} is 4 bytes, \code{0e 00 00 00}, identical in all 1,169 objects examined, and the last instruction
  of every program; composed multi-GEMM programs remove intermediate ENDs (decoder; Apple's compiler emits; MM~\mm{15}).
\item \op{13483}, \code{06 00}, is the compiler's nop and pads bodies; nearly every kernel contains it
  (142,452 instances in one corpus) (MM~\mm{25.68}).
\item Per-kernel metadata slot 32 is absent for main programs of 30 instructions or fewer
  and present (valued by length and back edge) above, a law of Apple's compiler on this toolchain, not a device
  guarantee (\cref{sec:gpu-metadata-slots}). While this compiler could author only the slot-present layout it padded bodies with two-byte
  fillers to at least 31 instructions (MM~\mm{15}); both layouts are now generated (\code{agxforge/g17/mdgen.py} \code{slot32\_for}).
\item Code is position-independent. An identical body is bit-exact at stream positions 1–12, after 0–4,096 bytes of padding and
  at offsets 2–40,002 bytes: no setup or alignment sequence is needed (measured, MM~\mm{15}). Metal places the code at offset
  +0x2740 of the dispatch's first mapped 64 KiB region (measured, MM~\mm{25.147}).
\item A branch target is offset + displacement (decoder; \cref{sec:branches-back-edges} for semantics), landing on a real instruction
  boundary for 1,867 of 1,867 corpus branches (offset plus size: 12.9 percent). \op{462} and
  \op{458} are both 10 bytes; the skip's bytes 5–9 are constant \code{00 00 00 00 00}, the back
  edge's \code{00 ff c7 ff 7f}. A back edge authored with a zeroed tail decoded cleanly but was wrong in its whole second half, and
  the error showed up only on execution (\code{docs/agx3-oracle.md} sections 12–13). A G17 branch names no condition; the condition comes from
  execution-mask state (\cref{ch:control-flow-execution}).
\end{itemize}

\section{Apple's decoder}\label{sec:decoder}

The per-generation compiler plugins (\code{applegpu\_g13p} to \code{applegpu\_g18p} on macOS 26) implement 30 of
32 plugin entry points; the two stubs are assembly emission and object emission, so \code{-S} is unreachable for every AGX
target. \code{air-objdump} reports "no instruction printer" for agx2/agx3 (and "no disassembler" for agx1), which implies
that the MCDisassembler exists. In \code{GPUCompiler.framework/Versions/32023/Libraries/libLLVM.dylib} the agx3 \code{llvm::Target} has its
disassembler constructor (+0x80) set and its printer constructor (+0x88) null; Apple withheld only the printer
(\code{docs/agx3-oracle.md} sections 1–3).

The decoder is reached by writing a printer factory into +0x88; \code{LLVMCreateDisasm} then succeeds. \code{printInst} is vtable slot 4,
found by pointing all 32 slots at recording stubs. The \code{MCInst} lives in \code{LLVMDisasmInstruction}'s frame, so operands must
be copied inside \code{printInst}: reading them after return gave 69.13 percent invalid operand kinds; copying inside gave 0
invalid over 89,645 operands in 290 objects. \code{tools/agx3meta.c} dumps the instruction,
operand, register-class and register tables (strides 48, 6, 32 and 24). No privileges, entitlements or patched binaries
are involved (\code{docs/agx3-oracle.md} sections 4–8). The in-process form is \code{tools/libagx3dis.dylib}, behind
\code{agxforge.g17.model.decode}; both this compiler's read-back guards and the project's emulator decode through it, so neither
runs on another machine (MM~\mm{25.209}).

Several things break on an OS update (\code{docs/agx3-oracle.md} section 9). The library path contains \code{Versions/32023}; the
\code{llvm::Target} field indices and every table stride are build properties (\code{agx3meta} re-validates strides and refuses;
\code{agx3dis} checks the disassembler constructor is non-null); the vtable slot of \code{printInst} is not checked at run time,
and a reorder would silently capture nothing ("opcode 0 on every line"). The register, class and instruction tables are
committed snapshots held equal to the binary's output by test, so after an update the tests fail rather than the
snapshots silently going stale. The update to macOS 27 moved the numbering and those tests failed
(\cref{sec:decoder-renumbering}).

MM~\mm{25.63} reports that "our decode" consumes 9,556 of 9,556 vendor bodies exactly and
agrees with Apple's boundaries on 400 of 400 sampled bodies (67,787 instructions of the classes once disputed). That
decode is \code{agxforge.g17.model.decode}, which calls Apple's decoder (MM~\mm{25.209}), so the result shows the wrapper reproduces
Apple's framing; it is not an independent confirmation. Neither of the two independent attempts matches Apple's decoder. One is the framer
(\cref{sec:framing}). The other, a decoder learned from Apple-decoded programs, framed and identified 99.98 percent of 746,000
held-out instructions but matched only 83.2 percent of operand tokens, left 16.7 percent unlearned, and got 636 (0.085
percent) wrong without any refusal, on bits that never varied in training; it is paused (MM~\mm{25.209}). An independent
decoder would have to be derived from per-length field maps with an encode-then-compare self-check that turns every uncertain decode into a
refusal (MM~\mm{25.209}).

\section{Opcode renumbering in macOS 27}\label{sec:decoder-renumbering}

macOS 27.0.1 (26A434) ships a rebuilt AGX3 LLVM at the same path. The decoder still runs, but TableGen re-sorted its
tables: the instruction table grew from 17,796 to \val{renumber-new-opcodes} entries and the register table from \val{renumber-old-registers} to \val{renumber-new-registers} entries,
so \op{684} decodes as 706. The encoding did not move. All \val{renumber-corpus-instructions} instructions of the 6,594 corpus
programs decode at their recorded offsets and lengths, and each old opcode the corpus uses lands on exactly one new one.
After the update 344 of the 722 test modules failed, among them the snapshot tests.

\code{tools/g17renumber.py} generates a table through which \code{agx3dis}, \code{libagx3dis} and \code{agx3meta} print
the 25G83 ids, so every opcode and register number in this reference refers to that build. Registers and register
classes map by name, since their names survived Apple's name stripping. Opcodes map from \val{renumber-witnessed} encodings decoded under
both builds: the corpus, the form record's witnesses and their opcode-bit flips, and the full decodes retained beside
Apple-compiled programs. Between those anchors the two descriptor tables are aligned, keeping only the placements that
every maximum alignment agrees on; \val{renumber-aligned} opcodes map that way and \val{renumber-ambiguous} are left unmapped as ambiguous. Built from the
corpus anchors alone, the alignment placed 2,684 of 2,687 held-out witnessed opcodes correctly, none wrongly, and left
three unmapped. The new build also writes a two-bit tag into bits 29--30 of some immediates, mostly in memory and
tensor forms; \val{renumber-tags} fitted entries undo it per opcode and operand. At start-up each decoder decodes eight canary
instructions and refuses a build whose canaries match neither numbering. With the translation the corpus decodes as
recorded, every operand value in the form record reads back (41,246 immediates and 29,635 registers), and 1,603
retained full decodes match token for token (\code{isa/g17-agx3-renumber.json}).

The new decoder also differs in ways that are not numbering effects. It rejects the repair-walk witnesses of
\op{9994} and \op{9995} (no glossary entry) and of \op{9996}, so D1 is \val{d1} on
26A434 against 7,549 on 25G83; Apple's own encoding of \opn{9996}, compiled from a 64-bit \code{atomic\_max}, still
decodes. Of the single-bit flips in the form record, it rejects 38 that the 25G83 decoder read as another opcode and
reads 7 that it rejected as opcodes new in the 26A434 build; operand classes are unchanged. And three operand slots
carry the bit-29 tag in every corpus instance with no record of their 25G83 values: \op{10023} operand 1, \op{10095} operand 2, and
\opn{9996} operand 1. They are decoded without translation, and no statement in this reference rests on them.

\section{Decoder admission and execution}\label{sec:admission-decodes-executes}

Neither a decoder rejection nor an acceptance settles what executes (\code{docs/agx3-oracle.md} section 12):

\begin{itemize}
\item An authored subtraction the decoder rejected computed 0 - x on hardware (the rejection was correct).
\item byte9{[}2{]} of the 14-byte load is rejected by the decoder yet executes as a narrow load: 0x12345678 with it clear,
  0x00000078 with it set, both predicted (measured).
\item Six decoder-refused one-bit variants of \op{5106} are no-ops on the accumulator and eight change the result; no encoding
  property separates the two groups, and the decode side is exhausted (measured, MM~\mm{2}, MM~\mm{25.25}, MM~\mm{25.36}).
\item Apple's compiler emits a decoder-refused encoding: int8 MMA streams contain \hash{a700a518220ea182c00a}; clearing structural
  bit 71 alone gives a 10-byte \op{10384}. With bit 71 set, the other 79 single flips
  never yield an MMA, and bit 71's meaning is open (MM~\mm{25}, MM~\mm{25.71}). The project's decoder wrapper names this variant
  and resumes ten bytes on.
\item Names R128–R143 decode and do not execute (\cref{sec:namespace}). A zeroed \op{458} tail decodes and is wrong (\cref{sec:program-boundaries}).
\item MMA admission is gated on the CPU name passed to the decoder: g17s, g17g and g18 admit the same ten tensor MMA opcodes;
  the 122 other table entries in scheduling classes 170–175 are unreachable within three flips of the seeds, a bounded
  negative (decoder, MM~\mm{2}, MM~\mm{25.4/F4} "F4, narrowed", MM~\mm{25.13}).
\item Metal's pipeline creation checks no instruction bytes: an undecodable \code{0xff} fill and a removed end-of-program both
  build (\cref{sec:object-library-archive}).
\end{itemize}

A bit the decoder ignores and Apple's corpus writes both ways can still change the result, so free bits need an
execution check. \op{11482}/14 bit 11.6 returns 1, 1, 1, 0 set and 1, 1, 1, 1 clear,
three runs each (measured, MM~\mm{25.88}). Drifting free-bit verdicts were adopted only after a hardware check (13 forms, then
two single-instance forms whose 28 and 11 bits were executed), 252 drift bits that change no emitted byte were
deliberately left unadopted, and the compare-select's bit stays live (MM~\mm{25.58}). Nor does a clean decode show that a form
works: forty repair-walk opcodes in one class ignore both inputs and return a constant of the encoding
(\code{docs/g17-isa-cartography.md}).

In the encode-then-decode-back guard, the encoder starts from a measured reference with its variable bits cleared,
solves each requested operand, emits, decodes the result through Apple's decoder, and refuses unless opcode, length and
every requested value round-trip (MM~\mm{23.3}). The guard was added after the encoder had relied on field arithmetic
alone: the campaign's encoder accepted a byte sequence whose field arithmetic reached a zero residual for an \emph{odd} register
base, and silently rounded it to the aligned tuple, because the tuple field cannot express the low bit (MM~\mm{2}, MM~\mm{23.3}).
Later read-back gaps of the same kind were found and closed: \op{9749}/8 asked for
register id 441 and wrote 425 until operands were read back; \op{1015}/4 accepted 2147483680
and emitted 32 because operands with an opcode-level code table were skipped (MM~\mm{25.88}). Every memory entry point in
\code{agxforge/g17/memenc.py} performs the same check (MM~\mm{23.3}). The guard checks only that encoder and decoder agree; it does not check
arithmetic, layout, wait masks or hardware admission (MM~\mm{23.3}).

The project's own dispatch gate refuses by Apple's instruction flags (may load, may store, branch, call, return,
terminator; memory, atomic, texture) and refuses opcodes with no Apple witness unless explicitly overridden; 83 of the 86
opcodes this compiler can emit carry an Apple witness (this compiler, MM~\mm{25.147}).

\textbf{Not known.} Per-length field maps for most of D1 (D2 is \val{d2} of \val{d1}); the 0xa1 class-7 lengths; the meaning of
int8-MMA bit 71; why 36 narrow-accumulator table rows are refused although structurally like admitted forms; the
remaining meaning of \op{12674}'s residue bit byte1{[}6{]} (MM~\mm{25}, MM~\mm{25.6}, MM~\mm{25.32}, MM~\mm{25.58}, MM~\mm{25.71}).

\textbf{Evidence.} MM~\mm{0.4}, MM~\mm{2}, MM~\mm{6}, MM~\mm{15}, MM~\mm{19}, MM~\mm{23.2}, MM~\mm{23.3}, MM~\mm{24.2} (row R6), 24.8 (row T1), 25 (introduction), 25.1, 25.4, 25.6, 25.13, 25.24,
25.25, 25.26, 25.32, 25.36, 25.37, 25.51, 25.52, 25.53, 25.58, 25.62, 25.63, 25.68, 25.71, 25.78, 25.81, 25.87, 25.88,
25.120, 25.120.1, 25.121, 25.147, 25.148, 25.196, 25.209; \code{docs/agx3-oracle.md}, \code{docs/agx-isa-status.md},
\code{docs/g17-isa-reference.md}, \code{isa/g17-tensor-mma-fieldmap.json}, \code{isa/g17-agx3meta-instrs.txt}, \code{agxforge/g17/memenc.py},
\code{agxforge/g17/mdgen.py}.
